% Hand-written discussion fragment for the control-amplifier subsection.
% NOT generated by maestro2tex -- the tool only places it, via the "order" list in
% configs/anatop_tb_controlamp.json.
%
% CAREFUL. Every number quoted here comes from tab_ctrlamp_largesignal.tex and
% tab_ctrlamp_dcop.tex, and they are SLEW figures. The three benches in this run
% contain no load sweep, so this subsection cannot state a DC drive rating and
% deliberately does not. If a load-sweep bench is added later, put its numbers in
% a generated table rather than transcribing them into this paragraph.

\paragraph{Output drive is asymmetric, and that is the open question on this block.}
The two slew directions differ by a factor of four, and the reason is structural rather
than a matter of bias current. When the output slews up, the drive comes from the
cascoded PMOS mirror, which can pass several times its quiescent current as the input
pair steers the folding node; when it slews down, the only pull-down at the output is
the NMOS folding current source, so the sink current can never exceed what that source
is biased to carry, less whatever the input pair is still delivering. Sourcing is
limited by a mirror ratio; sinking is limited by a fixed bias current. Moving the trim
word changes both --- \texttt{BiasAdj} is inverted, so \emph{lower} codes raise the
reference current --- but it moves the quiescent current with them and does not change
the ratio.

This matters because the rev-2 potentiostat asks A1 to drive the counter electrode
directly and bipolarly, against a full-scale target of \(\pm\SI{33}{\micro\ampere}\) at
the \SI{35}{\kilo\ohm} transimpedance setting. Nothing in the small-signal results is
the obstacle: the gain, the bandwidth and the \SI{65}{\degree} phase margin are all
comfortable, and the amplifier draws only \SI{31.8}{\micro\ampere} doing it. The
question is entirely the output branch.

What this run does \emph{not} establish is how much current the amplifier can hold at
DC. All three benches here load the output with \SI{10}{\pico\farad} and nothing else,
so the currents in Table~\ref{tab:ctrlamp-largesignal} are slewing currents measured
with the input pair slammed past its common-mode range --- a generous upper bound on the
output stage, not a compliance figure. DC drive has to be measured by sweeping a load
current at the output and watching the tracking error, and that bench is not part of
this run. Until it is, treat the asymmetry above as the shape of the problem and not as
its magnitude, and do not read the sourcing figure as headroom against the
\(\pm\SI{33}{\micro\ampere}\) target. Settling this is the item that should gate
committing this amplifier to the four-channel layout.
